\lesson{3}{Sep 10 2021 Fri (06:53:22)}{Energy and Temperature}

Particles in motion can cause change due to their motion. For example:

\begin{definition}[Energy]
		\textbf{Energy} is the ability to cause change or the ability to do work.
	\textbf{Energy} is measured in the \textbf{SI Unit Joules}:

	\[ \textrm{Joule} = \frac{kg \times m^2}{s^^2} .\]
\end{definition}

\begin{definition}[Work]
	\textbf{Work} is the process of causing matter to move against an opposing
	force. When matter moves, it's position is changing. For that to happen, it
	requires \textbf{Energy}.
\end{definition}

All forms of energy can be broadly classified as either:

\begin{itemize}
	\item \textbf{Potential Energy}.
	\item \textbf{Kinetic Energy}.
\end{itemize}

\subsection*{Different Kinds of Energy}

\begin{definition}[Potential Energy]
	\textbf{Potential Energy} is the energy an object has because of it's current
	position or composition. For example, you have someone standing on a plane
	$25,000 ft$ above ground before he sky dives which gives in \textbf{Potential Energy}.

	\textbf{Potential Energy} is also known as \textbf{Stored Energy} because
	it's energy that is there, but hasn't been used. You can think about it as,
	winding up a toy car, then once you let it go, all of the energy comes out as
	\textbf{Kinetic Energy}.
\end{definition}

\begin{definition}[Kinetic Energy]
	\textbf{Kinetic Energy} is energy an object
	has due to its motion. Like I mentioned before, when you wind up a toy car
	and then let it go, that's an example of \textbf{Kinetic Energy}.
\end{definition}

\begin{definition}[Mechanical Energy]
	\textbf{Mechanical Energy} is the movement of objects from one place to
	another

	EXAMPLES:
	\begin{itemize}
		\item Turning wheels.
		\item Wind blowing.
	\end{itemize}
\end{definition}

\begin{definition}[Stored Mechanical Energy]
	\textbf{Stored Mechanical Energy} is the energy stored in objects by
	application of force.

	EXAMPLES:
	\begin{itemize}
		\item Compressed spring.
		\item Pulled back string of a bow before shooting the arrow.
	\end{itemize}
\end{definition}

\begin{definition}[Electrical Energy]
	\textbf{Electrical Energy} is the movement of electrical charge.

	EXAMPLES:
	\begin{itemize}
		\item Electricity running though wires.
		\item Lighting.
	\end{itemize}
\end{definition}

\begin{definition}[Electrical Potential Energy]
	\textbf{Electrical Potential Energy} is the stored energy of electric charges
	due to their position and interaction around other charges.

	EXAMPLES:
	\begin{itemize}
		\item Non-moving charges in circuits.
	\end{itemize}
\end{definition}

\begin{definition}[Radiant Energy]
	\textbf{Radiant Energy} is electromagnetic energy that travels in waves.

	EXAMPLES:
	\begin{itemize}
		\item X-Rays.
		\item Microwaves.
		\item Visible light.
		\item Radio waves.
	\end{itemize}
\end{definition}

\begin{definition}[Chemical Energy]
	\textbf{Chemical Energy} is stored in the bonds of atoms and molecules.

	EXAMPLES:
	\begin{itemize}
		\item Biomass.
		\item Natural gas.
		\item Propane.
	\end{itemize}
\end{definition}

\begin{definition}[Sound Energy]
	\textbf{Sound Energy} is the movement of energy through substances in
	longitudinal waves.  Sound is produced when a force causes an object to move
	or vibrate, which is when energy is transferred through the object in a wave.
\end{definition}

\begin{definition}[Nuclear Energy]
	\textbf{Nuclear Energy} is energy stored in the nucleus of an atom. This is
	the energy that holds atoms together and can be released in processes called
	fusion and fission.
\end{definition}

\begin{definition}[Thermal Energy]
	\textbf{Thermal Energy} also known as \textbf{heat}, \textbf{Thermal Energy}
	is the internal energy in a substance caused by vibration of atoms and
	molecules. \textbf{Geothermal Energy} is an example of \textbf{Thermal
	Energy}. Heat is also a byproduct of a lot of energy conversions.
\end{definition}

\begin{definition}[Gravitational Energy]
	\textbf{Gravitational Energy} is the \textbf{Potential Energy} of position.
	The rock resting at the top of a hill has \textbf{Gravitational Potential
	Energy}.
\end{definition}

Although there are many forms of energy, the \textbf{Law of Conservation of Energy} states that
\textit{\textbf{"Energy can be converted from one form to another, but it's not created or destroyed in ordinary physical and chemical processes."}}

